\documentclass[11pt,a4paper]{article}
% preamble.tex — Shared LaTeX preamble for all tcfs design documents
% Usage: % preamble.tex — Shared LaTeX preamble for all tcfs design documents
% Usage: % preamble.tex — Shared LaTeX preamble for all tcfs design documents
% Usage: \input{preamble} at top of each .tex file

% Encoding and fonts
\usepackage[T1]{fontenc}
\usepackage[utf8]{inputenc}
\usepackage{lmodern}
\usepackage{textcomp}

% Page geometry
\usepackage[margin=2.5cm]{geometry}

% Colors and hyperlinks
\usepackage[dvipsnames]{xcolor}
\usepackage[colorlinks=true,linkcolor=NavyBlue,urlcolor=NavyBlue,citecolor=NavyBlue]{hyperref}

% Code listings
\usepackage{listings}
\lstset{
  basicstyle=\ttfamily\small,
  breaklines=true,
  frame=single,
  backgroundcolor=\color{gray!10},
  xleftmargin=2em,
  framexleftmargin=1.5em,
  numbers=none,
  showstringspaces=false,
  columns=flexible,
  literate={←}{$\leftarrow$}1 {→}{$\rightarrow$}1 {⟶}{$\longrightarrow$}1,
}
\lstdefinelanguage{json}{
  string=[s]{"}{"},
  comment=[l]{//},
  morecomment=[s]{/*}{*/},
  literate=
    *{:}{{{\color{NavyBlue}:}}}{1}
    {,}{{{\color{NavyBlue},}}}{1}
    {\{}{{{\color{NavyBlue}\{}}}{1}
    {\}}{{{\color{NavyBlue}\}}}}{1}
    {[}{{{\color{NavyBlue}[}}}{1}
    {]}{{{\color{NavyBlue}]}}}{1},
}
\lstdefinelanguage{toml}{
  comment=[l]{\#},
  string=[b]",
  morestring=[b]',
}
\lstdefinelanguage{protobuf}{
  morekeywords={service,rpc,returns,message,string,uint32,uint64,int64,bool,bytes,repeated,optional,oneof,enum,map},
  morecomment=[l]{//},
  morestring=[b]",
}

% Tables
\usepackage{longtable}
\usepackage{booktabs}
\usepackage{array}
\usepackage{tabularx}

% Graphics
\usepackage{graphicx}
\usepackage{float}

% Better lists
\usepackage{enumitem}

% Header/footer
\usepackage{fancyhdr}
\pagestyle{fancy}
\fancyhf{}
\fancyfoot[C]{\thepage}
\renewcommand{\headrulewidth}{0.4pt}

% Tighter spacing
\setlength{\parindent}{0pt}
\setlength{\parskip}{6pt}

% Convenience macros
\newcommand{\tcfs}{\texttt{tcfs}}
\newcommand{\tcfsd}{\texttt{tcfsd}}
\newcommand{\code}[1]{\texttt{#1}}
\newcommand{\filepath}[1]{\texttt{#1}}
 at top of each .tex file

% Encoding and fonts
\usepackage[T1]{fontenc}
\usepackage[utf8]{inputenc}
\usepackage{lmodern}
\usepackage{textcomp}

% Page geometry
\usepackage[margin=2.5cm]{geometry}

% Colors and hyperlinks
\usepackage[dvipsnames]{xcolor}
\usepackage[colorlinks=true,linkcolor=NavyBlue,urlcolor=NavyBlue,citecolor=NavyBlue]{hyperref}

% Code listings
\usepackage{listings}
\lstset{
  basicstyle=\ttfamily\small,
  breaklines=true,
  frame=single,
  backgroundcolor=\color{gray!10},
  xleftmargin=2em,
  framexleftmargin=1.5em,
  numbers=none,
  showstringspaces=false,
  columns=flexible,
  literate={←}{$\leftarrow$}1 {→}{$\rightarrow$}1 {⟶}{$\longrightarrow$}1,
}
\lstdefinelanguage{json}{
  string=[s]{"}{"},
  comment=[l]{//},
  morecomment=[s]{/*}{*/},
  literate=
    *{:}{{{\color{NavyBlue}:}}}{1}
    {,}{{{\color{NavyBlue},}}}{1}
    {\{}{{{\color{NavyBlue}\{}}}{1}
    {\}}{{{\color{NavyBlue}\}}}}{1}
    {[}{{{\color{NavyBlue}[}}}{1}
    {]}{{{\color{NavyBlue}]}}}{1},
}
\lstdefinelanguage{toml}{
  comment=[l]{\#},
  string=[b]",
  morestring=[b]',
}
\lstdefinelanguage{protobuf}{
  morekeywords={service,rpc,returns,message,string,uint32,uint64,int64,bool,bytes,repeated,optional,oneof,enum,map},
  morecomment=[l]{//},
  morestring=[b]",
}

% Tables
\usepackage{longtable}
\usepackage{booktabs}
\usepackage{array}
\usepackage{tabularx}

% Graphics
\usepackage{graphicx}
\usepackage{float}

% Better lists
\usepackage{enumitem}

% Header/footer
\usepackage{fancyhdr}
\pagestyle{fancy}
\fancyhf{}
\fancyfoot[C]{\thepage}
\renewcommand{\headrulewidth}{0.4pt}

% Tighter spacing
\setlength{\parindent}{0pt}
\setlength{\parskip}{6pt}

% Convenience macros
\newcommand{\tcfs}{\texttt{tcfs}}
\newcommand{\tcfsd}{\texttt{tcfsd}}
\newcommand{\code}[1]{\texttt{#1}}
\newcommand{\filepath}[1]{\texttt{#1}}
 at top of each .tex file

% Encoding and fonts
\usepackage[T1]{fontenc}
\usepackage[utf8]{inputenc}
\usepackage{lmodern}
\usepackage{textcomp}

% Page geometry
\usepackage[margin=2.5cm]{geometry}

% Colors and hyperlinks
\usepackage[dvipsnames]{xcolor}
\usepackage[colorlinks=true,linkcolor=NavyBlue,urlcolor=NavyBlue,citecolor=NavyBlue]{hyperref}

% Code listings
\usepackage{listings}
\lstset{
  basicstyle=\ttfamily\small,
  breaklines=true,
  frame=single,
  backgroundcolor=\color{gray!10},
  xleftmargin=2em,
  framexleftmargin=1.5em,
  numbers=none,
  showstringspaces=false,
  columns=flexible,
  literate={←}{$\leftarrow$}1 {→}{$\rightarrow$}1 {⟶}{$\longrightarrow$}1,
}
\lstdefinelanguage{json}{
  string=[s]{"}{"},
  comment=[l]{//},
  morecomment=[s]{/*}{*/},
  literate=
    *{:}{{{\color{NavyBlue}:}}}{1}
    {,}{{{\color{NavyBlue},}}}{1}
    {\{}{{{\color{NavyBlue}\{}}}{1}
    {\}}{{{\color{NavyBlue}\}}}}{1}
    {[}{{{\color{NavyBlue}[}}}{1}
    {]}{{{\color{NavyBlue}]}}}{1},
}
\lstdefinelanguage{toml}{
  comment=[l]{\#},
  string=[b]",
  morestring=[b]',
}
\lstdefinelanguage{protobuf}{
  morekeywords={service,rpc,returns,message,string,uint32,uint64,int64,bool,bytes,repeated,optional,oneof,enum,map},
  morecomment=[l]{//},
  morestring=[b]",
}

% Tables
\usepackage{longtable}
\usepackage{booktabs}
\usepackage{array}
\usepackage{tabularx}

% Graphics
\usepackage{graphicx}
\usepackage{float}

% Better lists
\usepackage{enumitem}

% Header/footer
\usepackage{fancyhdr}
\pagestyle{fancy}
\fancyhf{}
\fancyfoot[C]{\thepage}
\renewcommand{\headrulewidth}{0.4pt}

% Tighter spacing
\setlength{\parindent}{0pt}
\setlength{\parskip}{6pt}

% Convenience macros
\newcommand{\tcfs}{\texttt{tcfs}}
\newcommand{\tcfsd}{\texttt{tcfsd}}
\newcommand{\code}[1]{\texttt{#1}}
\newcommand{\filepath}[1]{\texttt{#1}}


\fancyhead[L]{\small\textsc{tcfs}}
\fancyhead[R]{\small Architecture}

\title{tcfs Architecture}
\author{tcfs / TummyCrypt}
\date{February 2026}

\begin{document}
\maketitle
\tableofcontents
\newpage

% ─────────────────────────────────────────────────────────────────────────────
\section{Overview}

\tcfs{} (TummyCrypt FileSystem) is a FOSS, self-hosted, S3-first synchronization
system that replaces proprietary cloud storage clients (odrive, Dropbox, etc.)
with platform-specific filesystem surfaces that transparently hydrate files on
demand.

\subsection{Design Goals}

\begin{enumerate}
  \item \textbf{Self-hosted first} --- all data stays on user-controlled infrastructure (SeaweedFS).
  \item \textbf{On-demand hydration} --- remote files appear locally as clean
        names or platform placeholders and are downloaded only when accessed.
  \item \textbf{End-to-end encryption} --- client-side XChaCha20-Poly1305 encryption
        before upload; storage operator never sees plaintext.
  \item \textbf{Content-addressed deduplication} --- FastCDC chunking with BLAKE3
        hashing eliminates redundant storage.
  \item \textbf{Cross-platform} --- Linux (FUSE3), macOS/iOS (FileProvider), Windows (Cloud Files API).
  \item \textbf{Kubernetes-native backend} --- stateless sync workers scaled by KEDA,
        NATS JetStream for reliable task dispatch.
\end{enumerate}

% ─────────────────────────────────────────────────────────────────────────────
\section{System Architecture}

The system is divided into client-side and server-side components connected
via S3 API and NATS JetStream.

\subsection{Client Side}

The client runs on user machines and consists of:

\begin{description}
  \item[\tcfsd{} (daemon)] Long-running process exposing a gRPC service over a
    Unix domain socket. Manages Linux FUSE mounts, credential loading, sync state,
    and Prometheus metrics. Reports readiness via \code{sd\_notify(READY=1)}.
  \item[\tcfs{} (CLI)] Command-line interface for push, pull, mount, unmount,
    sync-status, and unsync operations.
  \item[\code{tcfs-tui} (TUI)] Interactive terminal dashboard built with
    ratatui~0.30, connecting to the daemon via gRPC. Five tabs: Dashboard,
    Config, Mounts, Secrets, Conflicts.
  \item[\code{tcfs-mcp} (MCP server)] Model Context Protocol server for AI
    agent integration. Seven non-resolution tools over stdio JSON-RPC transport.
\end{description}

\subsection{Server Side (Kubernetes Model)}

This is the design model. Live state differs: the lab backend (one SeaweedFS
deployment plus NATS) has been paused at zero replicas since 2026-09-29, the
sync worker is not part of the planned return, and tcfs will come back on a
new SeaweedFS deployment. See \texttt{docs/ops/current.md}.

\begin{description}
  \item[SeaweedFS] S3-compatible blob storage (master, volume and filer
    roles).
  \item[NATS JetStream] Reliable task queue with two streams:
    \code{SYNC\_TASKS} and \code{HYDRATION\_EVENTS}.
  \item[Sync workers] Stateless \tcfsd{} instances running in
    \code{--mode=worker}. Designed for horizontal KEDA scaling based on NATS
    queue depth; full rollout proof is still tracked separately.
  \item[Prometheus + Grafana] Observability stack monitoring throughput,
    queue depth, and FUSE latency.
\end{description}

% ─────────────────────────────────────────────────────────────────────────────
\section{Crate Map}

All Rust code lives in \filepath{crates/}. The workspace contains 18 product
crates plus an end-to-end test crate:

\begin{longtable}{llp{8cm}}
  \toprule
  \textbf{Crate} & \textbf{Type} & \textbf{Purpose} \\
  \midrule
  \endhead
  \code{tcfs-core}        & library & Shared types, config schema, protobuf defs \\
  \code{tcfs-crypto}      & library & XChaCha20-Poly1305, Argon2id KDF, HKDF, BIP-39 \\
  \code{tcfs-auth}        & library & TOTP, WebAuthn/FIDO2, session, and enrollment helpers \\
  \code{tcfs-secrets}     & library & SOPS/age/KDBX credential chain \\
  \code{tcfs-storage}     & library & OpenDAL abstraction, SeaweedFS native API \\
  \code{tcfs-chunks}      & library & FastCDC chunking, BLAKE3 hashing (its zstd module is unused) \\
  \code{tcfs-sync}        & library & Sync engine, JSON state cache, NATS consumers \\
  \code{tcfs-vfs}         & library & Shared virtual filesystem, disk cache, stubs, hydration \\
  \code{tcfs-fuse}        & library & Linux FUSE driver (fuse3 crate) \\
  \code{tcfs-nfs}         & library & NFS loopback mount backend \\
  \code{tcfs-cloudfilter} & library & Windows Cloud Files API (skeleton) \\
  \code{tcfs-sops}        & library & SOPS+age fleet secret propagation \\
  \code{tcfs-file-provider} & library & macOS/iOS FileProvider FFI (RFC 0002) \\
  \code{tcfs-dbus}        & library & Linux D-Bus desktop status integration \\
  \code{tcfsd}            & binary  & Daemon: gRPC, Linux FUSE, metrics, systemd \\
  \code{tcfs-cli}         & binary  & CLI: push, pull, mount, status, unsync \\
  \code{tcfs-tui}         & binary  & TUI: ratatui dashboard \\
  \code{tcfs-mcp}         & binary  & MCP server: 7 non-resolution tools, stdio transport \\
  \bottomrule
\end{longtable}

% ─────────────────────────────────────────────────────────────────────────────
\section{Physical Stub File Format}

Physical sync-root and compatibility paths can represent files as \code{.tc}
stubs. Mounted VFS views normally expose clean filenames and hydrate from remote
index entries; macOS FileProvider uses platform placeholders / APFS dataless
files. A physical \code{.tc} stub is a small sorted key/value text file:

\begin{lstlisting}
version https://tummycrypt.io/tcfs/v1
chunks 23
compressed 0
fetched 0
oid blake3:4d7a214614ab2935c943f9e0ff69d22eadbb8f32b1258daaa5e2ca24d17e239
origin seaweedfs://filer.example.com/bucket/path/to/file
size 94371840
\end{lstlisting}

Directory stubs use the \code{.tcf} extension for physical directory
dehydration/compatibility paths. Mounted directory listings normally come from
remote \filepath{\{prefix\}/index/...} entries.

% ─────────────────────────────────────────────────────────────────────────────
\section{Hydration Sequence}

When a remote-backed file is opened through the mounted VFS:

\begin{enumerate}
  \item The mount backend's \code{open()} handler intercepts the file access
        request.
  \item \tcfsd{} reads the index entry at
        \filepath{\{prefix\}/index/\{rel\_path\}} and fetches the manifest it
        names from \filepath{\{prefix\}/manifests/\{manifest\_object\_id\}}
        (direct, non-indexed uploads use
        \filepath{\{prefix\}/manifests/\{file\_hash\}}).
  \item Chunks are fetched in parallel from
        \filepath{\{prefix\}/chunks/\{chunk\_hash\}}.
  \item Each chunk is verified (BLAKE3) and decrypted (XChaCha20-Poly1305).
        Chunks are not compressed.
  \item Chunks are concatenated in order to reconstruct the original file.
  \item The file is served to the caller and cached for future reads.
  \item Unsync/dehydrate evicts cached content or creates a physical
        \code{.tc} stub, depending on the surface.
\end{enumerate}

% ─────────────────────────────────────────────────────────────────────────────
\section{Credential Chain}

\tcfsd{} discovers age identities in order of precedence:

\begin{enumerate}
  \item \code{\$CREDENTIALS\_DIRECTORY/age-identity} --- systemd
        \code{LoadCredentialEncrypted}
  \item \code{\$SOPS\_AGE\_KEY\_FILE} --- path to an age key file
  \item \code{\$SOPS\_AGE\_KEY} --- literal age key content
  \item \code{secrets.age\_identity} --- path from the config file
  \item \filepath{\textasciitilde/.config/sops/age/keys.txt} --- default fallback
\end{enumerate}

Once loaded, the age identity decrypts SOPS-encrypted credential YAML files,
providing S3 access keys to the OpenDAL operator. An mtime watcher
auto-reloads credentials when the file changes on disk.

% ─────────────────────────────────────────────────────────────────────────────
\section{Phase Roadmap}

\begin{longtable}{clp{9cm}}
  \toprule
  \textbf{Phase} & \textbf{Status} & \textbf{Scope} \\
  \midrule
  \endhead
  0 & Complete & Repo foundation, SOPS migration, Rust workspace stubs \\
  1 & Complete & Core daemon + secrets + gRPC \\
  2 & Complete & Sync engine + chunking + NATS \\
  3 & Complete & FUSE driver + clean VFS hydration + physical stubs \\
  4 & Scaffolded; paused & K8s backend + HPA + OpenTofu deploy models; the lab backend is paused since 2026-09-29 and the worker is not part of the planned return \\
  5 & Complete & Release pipeline + packaging + docs site \\
  6 & Partial & macOS FileProvider lab lane + Windows CFAPI skeleton \\
  7 & Partial & Benchmark snapshot, fleet ops, Darwin RFC, production hardening \\
  \bottomrule
\end{longtable}

% ─────────────────────────────────────────────────────────────────────────────
\section{Infrastructure}

Host addresses, cluster identifiers and the live deployment inventory are
not published in this public repository. The lab's private inventory is
authoritative. Live backend state and the path back from the 2026-09-29
pause are in \texttt{docs/ops/current.md}.

Earlier deployment shapes (a bare-metal SeaweedFS cluster and a hosted
Kubernetes cluster with in-cluster SeaweedFS at
\code{seaweedfs.tcfs.svc.cluster.local:8333} and NATS at
\code{nats.tcfs.svc.cluster.local:4222}, image
\code{ghcr.io/jesssullivan/tcfsd}) are historical.

\end{document}
